% ===== 原文 前置部分（书名页 / 前言 ix–xiv / 引论 xi–xvi，PDF p.1–15，Lead） =====

% ---------- 书名页 ----------
\thispagestyle{empty}
\begin{center}
\vspace*{3cm}
{\Huge\bfseries 热带几何与镜像对称\par}
\vspace{1.2cm}
{\Large Mark Gross 著\par}
\vspace{0.6cm}
{\large 中文译注本\par}
\vspace{2cm}
{\small AMS CBMS Regional Conference Series in Mathematics, Volume 114\par}
\vspace{0.2cm}
{\small American Mathematical Society, Providence, Rhode Island\par}
\vspace{1.5cm}
{\small 原书：\textit{Tropical Geometry and Mirror Symmetry}, Mark Gross, 2011\par}
\vspace{0.3cm}
{\small 本译注本仅供学习研究使用；版权归原作者与 AMS 所有。\par}
\end{center}
\clearpage

% ---------- 目录 ----------
\tableofcontents
\clearpage

% ---------- 译者说明 ----------
\section*{译者说明}
\addcontentsline{toc}{section}{译者说明}
本译注本译自 Mark Gross 的 \textit{Tropical Geometry and Mirror Symmetry}（AMS 数学区域会议系列 114 卷，2011 年）。
原书基于作者 2008 年 12 月在堪萨斯州立大学 NSF-CBMS 会议上的十小时讲座，系统讲解热带几何、A/B 模型与对数几何，
并以 $\PP^2$ 为例完整走通 Mikhalkin 曲线计数与周期积分的 B-模型计算，最后概述作者与 Siebert 的 Gross--Siebert 纲领。

翻译遵循以下约定：
\begin{itemize}[leftmargin=1.6em,itemsep=1pt]
\item 术语首次出现时以“中文（English）”给出，之后统一用中文；人名、模型名与通行缩写（$\cN=2$、SYZ、GW 等）保留原文。
\item 公式编号、定理/定义/例/引理/命题/推论编号、图号一律与原文一致；原文的章内图号（如第 1 章 Figure 2）照旧。
\item 原文的笔误、省略的论证与可核对的易错点，以译者批注（\verb|% [Tutor Note: …]|）的形式写在源码中，排版时不显示。
\item 参考文献保留英文原貌；符号索引与总索引从略。
\end{itemize}

\clearpage

% ---------- 前言（原文 pp.ix–x）----------
\section*{前言}
\addcontentsline{toc}{section}{前言}

关于热带几何与镜像对称的 NSF-CBMS 会议于 2008 年 12 月 13 日至 17 日在堪萨斯州立大学举行。
会议的组织者是 Ricardo Castaño-Bernard、Yan Soibelman 与 Ilia Zharkov。会议期间，我作了十小时的讲座；
此外，M. Abouzaid、K.-W. Chan、C. Doran、K. Fukaya、I. Itenberg、L. Katzarkov、A. Mavlyutov、D. Morrison、
Y.-G. Oh、T. Pantev、B. Siebert 与 B. Young 也作了报告。

我的讲座意在给出一个长期研究纲领的概貌。这一纲领是我与 Bernd Siebert 正在推进的，目标是获得对镜像对称的根本性概念理解。
热带几何在这一纲领中自然地出现，因此在讲座中我采取了一种相当“非历史”的视角：从热带半环出发，发展热带几何，
解释 Mikhalkin 的热带曲线计数公式，并概述 Nishinou 与 Siebert 给出的证明；然后说明我近期将这一图景与镜像一侧联系起来的
工作；最后勾画我与 Siebert 近期工作的想法——如何从带奇点的仿射流形构造 Calabi--Yau 流形的退化。

本书紧密遵循讲座的结构，并补足了许多讲座中未能给出的细节。与讲座一样，本书只呈现一个仍在发展中的纲领的一个截面，
但我希望它对于愿意参与这一纲领的读者有所帮助。NSF 项目 DMS-0735319 为会议与本书提供了资助。

\vspace{0.8em}
\noindent\hfill Mark Gross，拉霍亚，2010 年

\clearpage

% ---------- 引论（原文 pp.xi–xvi）----------
\section*{引论}
\addcontentsline{toc}{section}{引论}

镜像对称的早期历史已被讲述过许多次；这里只作简要概括。故事始于 1985 年弦论中 Calabi--Yau 紧致化的引入 [11]。
其想法是：既然超弦理论要求十维时空，就要求时空（至少在局部上）取
\[
  \R^{1,3}\times X
\]
的形式，其中 $\R^{1,3}$ 是通常的 Minkowski 时空，$X$ 是一个很小的六维 Riemann 流形。
为使理论保持超弦理论的超对称性，就要求 $X$ 具有 $SU(3)$ 和乐，即 $X$ 是 Calabi--Yau 流形。
于是弦论进入了代数几何的领域：由 Yau 对 Calabi 猜想的证明 [113]，任何典范丛平凡的非奇异射影三维代数簇都带有 $SU(3)$ 和乐的度量。

这在弦论共同体中催生了一项产业：产生大量 Calabi--Yau 三维流形的例子并计算其不变量，其中最基本的是 Hodge 数 $h^{1,1}$ 与 $h^{1,2}$。
1989 年，从这项工作中浮现出一个相当令人惊讶的现象。Candelas、Lynker 与 Schimmrigk [12] 给出了一份加权射影空间中
Calabi--Yau 超曲面的列表，其中呈现出明显的对称性：若存在 Hodge 数为 $(h^{1,1},h^{1,2})$ 的 Calabi--Yau 三维流形，
则往往也存在 Hodge 数为 $(h^{1,2},h^{1,1})$ 的 Calabi--Yau 三维流形。几乎同时，Greene 与 Plesser [36] 在共形场论的
某些观察的引导下，研究了五次三维流形及其镜像伙伴。设 $X_\psi$ 是方程
\[
  x_0^5+\cdots+x_4^5-\psi x_0x_1x_2x_3x_4=0
\]
在 $\PP^4$ 中的解集，其中 $\psi\in\C$。则对大多数 $\psi$，$X_\psi$ 是非奇异五次三维流形，从而有 Hodge 数
\[
  h^{1,1}(X_\psi)=1,\qquad h^{1,2}(X_\psi)=101 .
\]
另一方面，群
\[
  G=\frac{\{(a_0,\dots,a_4)\mid a_i\in\mu_5,\ \prod_{i=0}^{4}a_i=1\}}
          {\{(a,a,a,a,a)\mid a\in\mu_5\}}
\]
通过
\[
  (x_0,\dots,x_4)\mapsto (a_0x_0,\dots,a_4x_4)
\]
对角地作用在 $\PP^4$ 上。这里 $\mu_5$ 是五次单位根群。该作用限制为 $X_\psi$ 上的作用，而商 $X_\psi/G$ 具有高度奇异性。
然而这些奇点可以通过一个双有理态射 $\check X_\psi\to X_\psi/G$ 消解，其中 $\check X_\psi$ 是新的 Calabi--Yau 三维流形，其 Hodge 数为
\[
  h^{1,1}(\check X_\psi)=101,\qquad h^{1,2}(\check X_\psi)=1 .
\]
这些例子在当时已令数学家惊讶，因为那时人们知道的 Euler 示性数为正的 Calabi--Yau 三维流形例子还很少
（Euler 示性数等于 $2(h^{1,1}-h^{1,2})$）。

更为惊人的是 Candelas、de la Ossa、Green 与 Parkes [10] 的结果。在弦论与路径积分计算的引导下，Candelas 等人猜测：
对以 $\psi$ 为参数的族 $\check X_\psi$ 作某些周期计算，可以得到五次三维流形上有理曲线数目的预言。他们完成了这些计算，
所得结果与已知的直到 3 次的有理曲线数目一致。这里略去这些计算的细节，因为它们在许多文献中已有阐述，例如见 [43]。
这一致性在数学界引起了极大惊讶，因为随着次数增加，这些数目变得越来越难以计算。直线数 2875 在 19 世纪就已知道；
二次曲线数 609250 直到 1986 年才由 Sheldon Katz [66] 算出；扭曲三次曲线数 317206375 直到 1990 年才由 Ellingsrud 与 Strømme [22] 算出。

在数学史中，物理一直是数学问题与数学现象的重要来源。由物理产生的一些有趣的数学往往只是“一次性的”——比如一个有趣而意外的公式，
一旦在数学上得到验证便失去了意义。物理的其他贡献则催生了强有力的新结构与新理论，并持续给出有趣而激动人心的新结果。
我愿相信镜像对称属于后者。
Candelas 等人提出的猜想以及相关工作，引导人们研究 Gromov--Witten 不变量（确切地定义什么是“有理曲线的数目”）与量子上同调
（用 Gromov--Witten 不变量对通常的上积进行形变的一种方式）。这至今仍是活跃的研究领域；到 1996 年，理论已发展到足以
使 Givental [34]、Lian--Liu--Yau [75] 以及后来的其他人证明五次三维流形的镜像对称公式，而且证明日趋简化。

关于镜像对称，Batyrev [6] 与 Batyrev--Borisov [7] 给出了非常一般的构造：在环面簇的完全交中出现的 Calabi--Yau 流形镜像对。
1994 年，Maxim Kontsevich [68] 提出了他的基本同调镜像对称猜想，这是用范畴论解释 Calabi--Yau 流形与其镜像之间关系的深刻努力。
1996 年，Strominger、Yau 与 Zaslow 提出了一个猜想 [108]，现称 SYZ 猜想，它暗示镜像对之间有一种更具体的几何关系：
镜像对应当带有对偶的特殊拉格朗日纤维化。这给出了 Calabi--Yau 流形与其镜像之间非常明确的关系；我本人 [37, 38, 39] 与
Wei-Dong Ruan [97, 98, 99] 在这一方向上的早期工作表明，该猜想在拓扑层面是成立的。然而至今，证明完整强度的 SYZ 猜想所需的
分析问题仍然难以逾越。此外，即使 SYZ 猜想得到证明固然极为重要，仅有一个证明并不能解释镜像对称更精细的方面。
尽管如此，SYZ 猜想启发了若干视角，它们看来正在为镜像对称提供新的洞见；尤其是 Kontsevich 与 Soibelman 在 [69, 70] 中
开创的刚性解析纲领，以及我与 Siebert 用对数几何发展的纲领 [47, 48, 51, 49]。

这些由 SYZ 猜想生长出来的想法聚焦于 SYZ 纤维化的底空间；尽管我们并不知道 SYZ 纤维化是否存在，但我们对其底空间的样子有很好的猜测。
特别地，它们应当是\emph{仿射流形}，即带有图册、其转移映射为仿射线性变换的实流形。一般而言，这些流形有一个奇异轨迹，
即不带有这种仿射结构的子集。写出这类流形的例子并不困难，我们期待它们对应例如环面簇中的超曲面。更确切地说：

\begin{definition}
\emph{仿射流形} $B$ 是一个实流形，带有一族坐标卡
\[
  \{\psi_i:U_i\to\R^n\},
\]
使得 $\psi_i\circ\psi_j^{-1}\in\operatorname{Aff}(\R^n)$，即 $\R^n$ 的仿射线性群。若
$\psi_i\circ\psi_j^{-1}\in\R^n\rtimes GL_n(\Z)\subseteq\operatorname{Aff}(\R^n)$（相应地，
$\psi_i\circ\psi_j^{-1}\in\operatorname{Aff}(\Z^n)$，即 $\Z^n$ 的仿射线性群），则称 $B$ 是\emph{热带的}
（相应地，\emph{整的}）。
\end{definition}

在热带情形，每个坐标变换的线性部分为整的；在整的情形，平移部分与线性部分都是整的。

给定热带流形 $B$，我们有一族格 $\Lambda\subseteq TB$，局部地由 $\partial/\partial y_1,\dots,\partial/\partial y_n$ 生成，
其中 $y_1,\dots,y_n$ 是仿射坐标。转移映射的条件保证这是良定义的。对偶地，我们有由 $\dd y_1,\dots,\dd y_n$ 生成的
一族格 $\check\Lambda\subseteq T^*B$，于是得到两个环面丛
\[
  f:X(B)\to B,\qquad \check f:\check X(B)\to B,
\]
其中
\[
  X(B)=TB/\Lambda,\qquad \check X(B)=T^*B/\Lambda .
\]
现在 $X(B)$ 带有自然的复结构。$\Lambda$ 的截面是 $TB$ 上某个联络的平坦截面，而该联络的水平与竖直切空间典范同构。
于是我们可以写出一个几乎复结构 $J$，它交换这两个空间，并带适当的符号变换使得 $J^2=-\id$。容易看出，
$TB$ 上的这个几乎复结构是可积的，并下降为 $X(B)$ 上的复结构。
另一方面，$T^*B$ 带有典范辛形式，它下降为 $\check X(B)$ 上的辛形式，因此 $\check X(B)$ 典则地是一个辛流形。

我们可以把 $X(B)$ 与 $\check X(B)$ 看作构成一个镜像对；这是 SYZ 猜想的一个简单版本。不过在这样简单的设定下，
有趣的紧例子很少：在 Kähler 情形，只限于 $B=\R^n/\Gamma$（$\Gamma$ 为格）这一可能（见 [15]）。
尽管如此，我们可以把这个简单情形当作动机，提出一些基本问题：
\begin{enumerate}[leftmargin=1.6em]
\item $B$ 上的哪些几何结构对应 $X(B)$ 与 $\check X(B)$ 上人们关心的几何结构？
\item 若要得到更有趣的例子，我们需要允许 $B$ 有奇点，即在开集 $B_0\subseteq B$ 上有热带仿射结构，而 $B\setminus B_0$ 相对较小
（例如余维至少为 2）。此时应当如何处理？
\end{enumerate}

到 2000 年，该领域的许多研究者已经清楚：$X(B)$ 中的全纯曲线应当对应 $B$ 中某种分片线性图。
Kontsevich 提出，也许能够通过数这些图来真正实现曲线计数。2002 年，Mikhalkin [79, 80] 宣布这确实是可行的，
并引入且证明了环面曲面的曲线计数公式。这是人们第一次真正能用这些分片线性图计算不变量。出于将在第 1 章解释的历史原因，
Mikhalkin 把这些分片线性图称为“热带曲线”，从而把这个词引入了该领域。

这就引出了如下图景。镜像对称涉及两类不同几何之间的关系，通常分别称为 A-模型与 B-模型。A-模型涉及辛几何，
这是讨论 Gromov--Witten 不变量之类对象的自然范畴；B-模型涉及复几何，可以在其中讨论周期积分之类对象。
这引导我们得到镜像对称的如下概念框架：
\begin{center}
\begin{tikzpicture}[scale=1.0,transform shape]
  \node (A) at (0,0) {A-模型};
  \node (B) at (4,0) {B-模型};
  \node (T) at (2,-1.8) {热带几何};
  \draw (A)--(B)--(T)--(A);
\end{tikzpicture}
\end{center}
这里我们希望这样解释镜像对称：在 $B$ 上辨认出我们称之为热带结构的结构，这些结构可以在 A-模型与 B-模型中解释为几何结构。
不过它们在 A-模型与 B-模型中的解释应当不同，即互为镜像；于是这些结构由同一热带结构给出这一事实，就为镜像对称提供了概念性的解释。
对于镜像对称最著名的方面，即有理曲线的枚举，我们希望 $B$ 上的热带曲线对应 A-模型中的（伪）全纯曲线，以及 B-模型中周期计算的修正项。

我与 Siebert 的纲领的主要想法是：通过另一个世界——对数几何的世界——来理解如何在热带世界与 A-模型、B-模型之间转换。
可以把对数几何看成热带几何与经典几何之间的中途：
\begin{center}
\begin{tikzpicture}[scale=1.0,transform shape]
  \node (A) at (0,0) {A-模型};
  \node (B) at (4,0) {B-模型};
  \node (L) at (2,-1.2) {对数几何};
  \node (T) at (2,-2.6) {热带几何};
  \draw (A)--(B);
  \draw (A)--(L)--(B);
  \draw (L)--(T) node[midway,right] {\small 几何};
\end{tikzpicture}
\end{center}

由于这一与 Siebert 合作的纲领仍在进行、仍有许多工作要做，我在堪萨斯州曼哈顿举行的 CBMS 区域会议上的讲座，
意在给出该纲领当前状态的截面。本书紧密遵循这些讲座的提纲。基本目标有三。
第一，我希望能至少在特殊情形下，把上图所提示的所有世界都明确解释清楚：热带世界、A-模型与 B-模型的“经典”世界，以及对数几何。
第二，我想解释一个非常具体的例子，其中 A-模型与 B-模型的完整图景都已被做出来，这就是 $\PP^2$ 的情形。
对 A-模型，曲线计数的结果是 Mikhalkin 的，这里我将给出他结果的一个证明；这个证明改编自 Nishinou--Siebert [86] 的更一般结果，
因为那种做法更符合本纲领的理念。对 B-模型，我将解释我自己近期的工作 [42]，它说明周期积分如何提取热带信息。
第三，我想综述我与 Siebert 在 Calabi--Yau 情形得到的一些结果，勾画这一进路为何有望给出镜像对称的证明。
对 $\PP^2$ 我给出完整细节；第三部分则更像是一份阅读原始论文的指南，因为那些论文不幸相当长且技术性很强。我希望至少能传达这一进路的直观。

在所有这些内容中，我将采取非常“非历史”的进路：从热带几何的基础出发，逆向展开，说明对热带几何的研究如何自然地引出那些最初在镜像对称研究中出现的概念。
某种意义上这或许是自然的。用 Witten 关于弦论的说法来转述：镜像对称常常像是 21 世纪的数学掉进了 20 世纪。
它最初在弦论中的发现，代表了该理论中较为困难的一些方面。哪怕只是解释 Candelas 等人所做的计算，也会占去一门课程相当大的部分；
而为定义与计算 Gromov--Witten 不变量所建立的理论则更为复杂。另一方面，如今看来构成镜像对称底层的几何——热带几何——却非常简单，理解它不需要特别的背景。
因此从最简单的起点展开讨论是合理的。

本书的先修要求包括：熟悉 Hartshorne 教科书 [57] 程度的代数几何，以及一些基本的微分几何。此外，熟悉环面几何将非常有帮助；
本书会回顾环面几何的许多基本事实，但一定程度的先前经验仍然有用。关于环面几何更深入的论述，我推荐 Fulton 的讲义 [27]。
第 3 章中我们还会用到平展拓扑（étale topology）中的层，可参阅 [83] 第 II 章；不过这一用法对本书的大部分讨论并不关键。

我要感谢许多人。首先感谢 Ricardo Castaño-Bernard、Yan Soibelman 与 Ilia Zharkov 在堪萨斯州立大学组织 NSF-CBMS 会议。
其次感谢 Bernd Siebert：本书中的进路源于我们共同的合作。我还要感谢许多回答问题并对书稿提出意见的人，包括 Sean Keel、
M. Brandon Meredith、Rahul Pandharipande、D. Peter Overholser、Daniel Schultheis 与 Katharine Shultis。
本书的部分内容写于访问牛津期间，感谢 Philip Candelas 的款待。本书最后部分写于 2009 年秋在 MSRI 期间，感谢 MSRI 通过 Simons 教授职位提供的资助。
感谢 Lori Lejeune 提供第 1 章 Figure 17、18、19 与第 6 章 Figure 7、8 的图像文件。最后但绝非最不重要的，感谢 Arthur Greenspoon，
他慷慨地提出为本书校对。

\para{约定。} 在本书中，$k$ 始终表示特征为零的代数闭域；$\N$ 表示自然数集 $\{0,1,2,\dots\}$。
